% GENERATED FILE. Edit canonical lessons, then run npm run generate:books.
% canonical-source-hash: fnv1a64:8edb43a6b5372105
% canonical-lessons: MR-C506-lipik, MR-C506-sipai, MR-C506-sikka, MR-C506-kagadpatra, MR-C506-channi-karne

\chapter{Office Life}
\label{ch:mr-office-life}
\hlchaptermodality{506}

\begin{chapteropening}
\textbf{By the end of this chapter:} \emph{I can talk about clerks, papers and stamps in an office.}

Complete the last lesson of chapter 506: I can talk about clerks, papers and stamps in an office.
\end{chapteropening}

\section[\texorpdfstring{\mr{लिपिक} (lipik)}{लिपिक (lipik)}]{\mr{लिपिक} (lipik) — a clerk}

\label{lesson:MR-C506-lipik}



\begin{quote}
Before the new one: say the Marathi for a weaver, then the Marathi for to set up.
\end{quote}

\subsection*{You'll want to know: \mr{लिपिक}}
\textbf{\mr{लिपिक}} — \emph{lipik} — \textquotedblleft{}a clerk\textquotedblright{}.

\begin{quote}
\mr{लिपिकाने} \mr{माझा} \mr{अर्ज} \mr{वाचला}.
\end{quote}

\begin{quote}
\emph{lipikāne mājhā arja vāclā.}
\end{quote}

\begin{quote}
The clerk read my application.
\end{quote}

Masculine.

\begin{cousinweb}
\textbf{\mr{लिपिक}} (\emph{lipik}) is a tatsama, taken from Sanskrit \emph{lipika}, \textquotedblleft{}a scribe, a clerk\textquotedblright{} (Monier-Williams), built on \emph{lipi}, \textquotedblleft{}writing\textquotedblright{}.
\end{cousinweb}

\begin{grammarlens}[title={What you've built}]
One more word for describing what you lived and felt.
\end{grammarlens}

\subsection*{Your turn}
\begin{itemize}
\raggedright
  \item \emph{Say it:} \emph{lipik}
  \item \emph{Say it:} \emph{lipik}, once more
  \item \emph{Say it:} \emph{lipikāne mājhā arja vāclā.} — The clerk read my application.
  \item \emph{Recall:} say the Marathi for a weaver, then the Marathi for to set up, then say \emph{lipik} again
\end{itemize}

\begin{tcolorbox}[breakable,colback=teal!4,colframe=teal!35!black,title={Before you move on}]
What does \mr{लिपिक} mean? (\textquotedblleft{}A clerk\textquotedblright{}.) Say it once more.
\end{tcolorbox}
\section[\texorpdfstring{\mr{शिपाई} (śipāī)}{शिपाई (śipāī)}]{\mr{शिपाई} (śipāī) — an office attendant; a peon}

\label{lesson:MR-C506-sipai}



\begin{quote}
Before the new one: say the Marathi for to set up, then the Marathi for a clerk.
\end{quote}

\subsection*{You'll want to know: \mr{शिपाई}}
\textbf{\mr{शिपाई}} — \emph{śipāī} — \textquotedblleft{}an office attendant; a peon\textquotedblright{}.

\begin{quote}
\mr{शिपायाने} \mr{सगळ्यांना} \mr{चहा} \mr{आणून} \mr{दिला}.
\end{quote}

\begin{quote}
\emph{śipāyāne sagḷyānnā cahā āṇūn dilā.}
\end{quote}

\begin{quote}
The office attendant brought everyone tea.
\end{quote}

Masculine.

The oblique is \mr{शिपाया}-: \mr{शिपायाने}, \mr{शिपायाला}.

\begin{cousinweb}
\textbf{\mr{शिपाई}} (\emph{śipāī}) is a loan from Persian \emph{sipāhī}, \textquotedblleft{}soldier\textquotedblright{}; in a Marathi office it names the attendant who carries files and tea.
\end{cousinweb}

\begin{grammarlens}[title={What you've built}]
One more word for describing what you lived and felt.
\end{grammarlens}

\subsection*{Your turn}
\begin{itemize}
\raggedright
  \item \emph{Say it:} \emph{śipāī}
  \item \emph{Say it:} \emph{śipāī}, once more
  \item \emph{Say it:} \emph{śipāyāne sagḷyānnā cahā āṇūn dilā.} — The office attendant brought everyone tea.
  \item \emph{Recall:} say the Marathi for to set up, then the Marathi for a clerk, then say \emph{śipāī} again
\end{itemize}

\begin{tcolorbox}[breakable,colback=teal!4,colframe=teal!35!black,title={Before you move on}]
What does \mr{शिपाई} mean? (\textquotedblleft{}An office attendant; a peon\textquotedblright{}.) Say it once more.
\end{tcolorbox}
\section[\texorpdfstring{\mr{शिक्का} (śikkā)}{शिक्का (śikkā)}]{\mr{शिक्का} (śikkā) — a rubber stamp; a seal}

\label{lesson:MR-C506-sikka}



\begin{quote}
Before the new one: say the Marathi for a clerk, then the Marathi for an office attendant.
\end{quote}

\subsection*{You'll want to know: \mr{शिक्का}}
\textbf{\mr{शिक्का}} — \emph{śikkā} — \textquotedblleft{}a rubber stamp; a seal\textquotedblright{}.

\begin{quote}
\mr{लिपिकाने} \mr{कागदावर} \mr{शिक्का} \mr{मारला}.
\end{quote}

\begin{quote}
\emph{lipikāne kāgdāvar śikkā mārlā.}
\end{quote}

\begin{quote}
The clerk stamped the paper.
\end{quote}

Masculine.

\mr{शिक्का} \mr{मारणे}, to stamp: the stamp is struck.

\begin{cousinweb}
\textbf{\mr{शिक्का}} (\emph{śikkā}) is a loan from Persian \emph{sikka}, \textquotedblleft{}a stamp, a coin die\textquotedblright{}, usually traced to Arabic \emph{sikka}.
\end{cousinweb}

\begin{grammarlens}[title={What you've built}]
One more word for describing what you lived and felt.
\end{grammarlens}

\subsection*{Your turn}
\begin{itemize}
\raggedright
  \item \emph{Say it:} \emph{śikkā}
  \item \emph{Say it:} \emph{śikkā}, once more
  \item \emph{Say it:} \emph{lipikāne kāgdāvar śikkā mārlā.} — The clerk stamped the paper.
  \item \emph{Recall:} say the Marathi for a clerk, then the Marathi for an office attendant, then say \emph{śikkā} again
\end{itemize}

\begin{tcolorbox}[breakable,colback=teal!4,colframe=teal!35!black,title={Before you move on}]
What does \mr{शिक्का} mean? (\textquotedblleft{}A rubber stamp; a seal\textquotedblright{}.) Say it once more.
\end{tcolorbox}
\section[\texorpdfstring{\mr{कागदपत्र} (kāgadpatra)}{कागदपत्र (kāgadpatra)}]{\mr{कागदपत्र} (kāgadpatra) — a document; papers}

\label{lesson:MR-C506-kagadpatra}



\begin{quote}
Before the new one: say the Marathi for an office attendant, then the Marathi for a rubber stamp.
\end{quote}

\subsection*{You'll want to know: \mr{कागदपत्र}}
\textbf{\mr{कागदपत्र}} — \emph{kāgadpatra} — \textquotedblleft{}a document; papers\textquotedblright{}.

\begin{quote}
\mr{मीराने} \mr{सगळी} \mr{कागदपत्रं} \mr{एका} \mr{पिशवीत} \mr{ठेवली}.
\end{quote}

\begin{quote}
\emph{mīrāne sagḷī kāgadpatraṁ ekā piśvīt ṭhevlī.}
\end{quote}

\begin{quote}
Meera kept all the papers in one bag.
\end{quote}

Neuter. Plural: \mr{कागदपत्रं}, \emph{kāgadpatraṁ}.

Mostly in the plural, \mr{कागदपत्रं}, the papers for a task.

\begin{cousinweb}
\textbf{\mr{कागदपत्र}} (\emph{kāgadpatra}) is a Marathi compound of \mr{कागद}, \textquotedblleft{}paper\textquotedblright{}, a loan from Persian \emph{kāghaz}, and \mr{पत्र}, \textquotedblleft{}a leaf, a letter\textquotedblright{}, from Sanskrit \emph{patra} (Monier-Williams).
\end{cousinweb}

\begin{grammarlens}[title={What you've built}]
One more word for describing what you lived and felt.
\end{grammarlens}

\subsection*{Your turn}
\begin{itemize}
\raggedright
  \item \emph{Say it:} \emph{kāgadpatra}
  \item \emph{Say it:} \emph{kāgadpatra}, once more
  \item \emph{Say it:} \emph{mīrāne sagḷī kāgadpatraṁ ekā piśvīt ṭhevlī.} — Meera kept all the papers in one bag.
  \item \emph{Recall:} say the Marathi for an office attendant, then the Marathi for a rubber stamp, then say \emph{kāgadpatra} again
\end{itemize}

\begin{tcolorbox}[breakable,colback=teal!4,colframe=teal!35!black,title={Before you move on}]
What does \mr{कागदपत्र} mean? (\textquotedblleft{}A document; papers\textquotedblright{}.) Say it once more.
\end{tcolorbox}
\section[\texorpdfstring{\mr{छाननी} \mr{करणे} (chānnī karṇe)}{छाननी करणे (chānnī karṇe)}]{\mr{छाननी} \mr{करणे} (chānnī karṇe) — to scrutinize}

\label{lesson:MR-C506-channi-karne}



\begin{quote}
Before the new one: say the Marathi for a rubber stamp, then the Marathi for a document.
\end{quote}

\subsection*{You'll want to know: \mr{छाननी} \mr{करणे}}
\textbf{\mr{छाननी} \mr{करणे}} — \emph{chānnī karṇe} — \textquotedblleft{}to scrutinize\textquotedblright{}.

\begin{quote}
\mr{लिपिकाने} \mr{सगळ्या} \mr{कागदपत्रांची} \mr{छाननी} \mr{केली}.
\end{quote}

\begin{quote}
\emph{lipikāne sagḷyā kāgadpatrāñcī chānnī kelī.}
\end{quote}

\begin{quote}
The clerk scrutinized all the papers.
\end{quote}

Past, with -\mr{ने} (he …): \mr{त्याने} \mr{अर्जाची} \mr{छाननी} \mr{केली}, \emph{tyāne arjācī chānnī kelī}.

What is checked takes -\mr{ची}: \mr{अर्जाची} \mr{छाननी} \mr{करणे}.

\begin{cousinweb}
\textbf{\mr{छाननी} \mr{करणे}} (\emph{chānnī karṇe}) is \mr{छाननी}, \textquotedblleft{}close sifting, scrutiny\textquotedblright{}, with \mr{करणे}; the origin of \mr{छाननी} is not settled.
\end{cousinweb}

\begin{grammarlens}[title={What you've built}]
One more word for describing what you lived and felt.
\end{grammarlens}

\subsection*{Your turn}
\begin{itemize}
\raggedright
  \item \emph{Say it:} \emph{chānnī karṇe}
  \item \emph{Say it:} \emph{chānnī karṇe}, once more
  \item \emph{Say it:} \emph{lipikāne sagḷyā kāgadpatrāñcī chānnī kelī.} — The clerk scrutinized all the papers.
  \item \emph{Recall:} say the Marathi for a rubber stamp, then the Marathi for a document, then say \emph{chānnī karṇe} again
\end{itemize}

\begin{tcolorbox}[breakable,colback=teal!4,colframe=teal!35!black,title={Before you move on}]
What does \mr{छाननी} \mr{करणे} mean? (\textquotedblleft{}To scrutinize\textquotedblright{}.) Say it once more.
\end{tcolorbox}
